%% ma: L10-B01-0019
%% lop: 10
%% chuong: 1
%% bai: 1
%% dang: 10.1.2
%% loai: TLN
%% muc: TH
%% dap_an: "1"
%% nguon: "Ảnh giáo viên gửi 10/10/2026 (ThS. Nguyễn Hoàng Việt, Chương 1 Mệnh đề), A-Đề 1, Phần 3 Câu 20"
%% kien_thuc: "Bài 1 (mệnh đề, mệnh đề chứa biến)"
%% trang_thai: cho-duyet
%% ly_do: "Dạng toán mới đề xuất 10.1.2: xét tính đúng sai của mệnh đề (số học, đại số, hình học)"
\begin{ex}
Trong các mệnh đề sau, có bao nhiêu mệnh đề sai?
\begin{enumerate}[label=\alph*)]
\item ``Số nguyên tố lớn hơn $2$ là số lẻ''.
\item ``Số tự nhiên có chữ số tận cùng là $0$ hoặc $5$ thì chia hết cho $5$''.
\item ``Bình phương tất cả các số nguyên đều chia hết cho $2$''.
\end{enumerate}
\shortans{1}
\loigiai{a) Đúng (số nguyên tố lớn hơn $2$ không chia hết cho $2$). b) Đúng (dấu hiệu chia hết cho $5$). c) Sai, chẳng hạn $1^2=1$ không chia hết cho $2$. Có $1$ mệnh đề sai.}
\end{ex}
